\subsection{The Unique Chain (Excerpt)}

The complete chain from $D = 58$ to $D = 8$:

\begin{center}
\small
\begin{tabular}{@{}clccc@{}}
\toprule
$n$ & \textbf{Orbit Label} & $D_n$ & $\dim(\mathcal{O})$ & $h$\\
\midrule
1 & $A_1$ & 58 & 190 & 2\\
2 & $2A_1$ & 56 & 192 & 4\\
3 & $3A_1$ & 54 & 194 & 6\\
$\vdots$ & $\vdots$ & $\vdots$ & $\vdots$ & $\vdots$\\
12 & (EW anchor) & 36 & 212 & 24\\
15 & (Hadronic) & 30 & 218 & 30\\
$\vdots$ & $\vdots$ & $\vdots$ & $\vdots$ & $\vdots$\\
26 & $E_8(a_1)$ & 8 & 240 & 58\\
\bottomrule
\end{tabular}
\end{center}

\textbf{Full data:} Available in machine-readable format at the replication repository 
(see Section~\ref{app:repro}).

\subsection{Verification Algorithm}

The uniqueness is verified by:
\begin{enumerate}[leftmargin=1.5em, topsep=2pt, itemsep=1pt]
    \item Constructing the complete Hasse diagram of $\mathcal{N}(\mathfrak{e}_8)$ 
        from standard tables.
    \item Computing all cover relations with $\Delta D = 2$.
    \item Running dynamic programming to find \emph{all} paths from $D = 58$ to $D = 8$.
    \item Applying the cost function $\mathcal{S} = \sum_n (\Delta^2 \ln D_n)^2$ 
        (second-difference smoothness).
    \item Verifying that the minimum-cost path is unique (cost gap to second-best $= 0$ 
        indicates tie-breaker was not needed for the optimal path).
\end{enumerate}

%=============================================================================
\section{Boundary Trace Derivation}
\label{app:boundary}
%=============================================================================

This appendix derives the block index formula $k_B = \frac{3}{2} I_1(B)$ from 
the APS index theorem.

\subsection{Setup}

Consider a 6D manifold $M^6$ with boundary $\partial M^6 = \Sigma^5$. For a 
Weyl fermion coupled to a $U(1)$ gauge field $A$, the APS index theorem gives:
\begin{equation}
    \text{ind}(\not{\!D}) = \int_{M^6} \text{ch}(F) \wedge \hat{A}(M^6) 
    - \frac{1}{2}\eta(\not{\!D}_\Sigma)
\end{equation}
where $\eta$ is the $\eta$-invariant of the boundary Dirac operator.

\subsection{Flat Connection on Boundary Cycles}

On the orientable double cover $\widetilde{M}$, the boundary decomposes as:
\begin{equation}
    \partial\widetilde{M} = C_1 \cup C_2 \cup C_T
\end{equation}
(two boundary cycles plus the twist seam).

For a \emph{flat} $U(1)$ connection on each $C_i$ with holonomy $\theta_i = \oint_{C_i} A$, 
the $\eta$-invariant contribution becomes:
\begin{equation}
    \eta_i \propto \sum_{\text{fermions}} q^2 \cdot \theta_i
\end{equation}
where $q$ is the $U(1)$ charge.

\subsection{\texorpdfstring{Spectral Trace Coefficient and $\bone$}{Spectral Trace Coefficient and b1}}
The regulated determinant of the hypercharge-coupled Dirac operator contains a universal
UV logarithmic term whose coefficient is the quadratic charge trace:
\begin{equation}
    \ln\det(\not{D}_Y) \supset \frac{1}{2\pi}\left(\frac{5}{3}\sum_f Y_f^2\right)
    \ln\frac{\Lambda}{\mu}\;\int F^2.
\end{equation}
Thus $\bone = \frac{5}{3}\sum_f Y_f^2 = 41/10$ appears as a \emph{spectral index coefficient},
independent of IR thresholds, matching its role in the CFE.

\subsection{\texorpdfstring{Derivation of $k_B$}{Derivation of kB}}

\textbf{Step 1:} Each boundary cycle contributes a factor proportional to $I_1 = \sum q^2$.

\textbf{Step 2:} Three cycles contribute additively: factor 3.

\textbf{Step 3:} The double cover counts each physical state once across two sheets. 
The effective weight is halved: factor $1/2$.

\textbf{Result:}
\begin{equation}
    k_B = 3 \times \frac{1}{2} \times I_1(B) = \frac{3}{2} I_1(B)
\end{equation}

\textbf{Consistency check:} For the electroweak block, 
$I_1^{\text{EW}} = \frac{41}{48}$ (SM hypercharge sum), giving 
$k_{\text{EW}} = \frac{3}{2} \times \frac{41}{48} = \frac{41}{32}$, 
as stated in Table~\ref{tab:blocks}.

%=============================================================================
\section{\texorpdfstring{Seam Gluing and the $\eta$-Invariant Normalization}{Seam Gluing and the eta-Invariant Normalization}}
\label{app:seam}
%=============================================================================

This appendix provides the $\eta$-invariant gluing derivation used in Axiom~\ref{ax:GB}.

\subsection{Cut Space and Operators}
Let $\widetilde{M}$ be the orientable double cover of the M\"obius fiber and cut it along
the seam $\Gamma$ to obtain a manifold-with-boundary $\widetilde{M}_\Gamma$ with
boundary components $(C_1,C_2,\Gamma_+,\Gamma_-)$. Let $\not{D}_\Gamma$ be the Dirac operator
on $\widetilde{M}_\Gamma$ with APS boundary conditions on $C_{1,2}$ and a unitary matching
operator $U_\Gamma$ identifying $\Gamma_+$ with $\Gamma_-$.
On the seam we take the tangential 1D Dirac operator
\begin{equation}
    D_\Gamma = i\,\frac{d}{d\theta}, \qquad \theta\in[0,2\pi).
\end{equation}

\subsection{\texorpdfstring{Gluing Formula for the $\eta$-Invariant}{Gluing Formula for the eta-Invariant}}
The APS gluing theorem expresses the $\eta$-invariant of the glued manifold in terms of
the cut manifold plus a discrete spectral-flow correction:
\begin{equation}
    \eta(\not{D}_{\widetilde{M}}) = \eta(\not{D}_\Gamma) + 2\,\mathrm{SF}(U_\Gamma),
\end{equation}
where $\mathrm{SF}(U_\Gamma)\in\mathbb{Z}$ is the spectral flow of the matching operator.
In TFPT bookkeeping this contributes a quantized seam term
\begin{equation}
    \Delta_\Gamma = 2\pi\,\mathrm{SF}(U_\Gamma)
\end{equation}
to the curvature normalization.

\subsection{M\"obius Matching and the Minimal Gluing Class}
For the M\"obius $\mathbb{Z}_2$ identification, the matching satisfies $U_\Gamma^2=\mathbf{1}$.
The minimal nontrivial class is modeled by the unitary phase
\begin{equation}
    U_\Gamma(\theta)=e^{i\theta}.
\end{equation}

\subsection{Spectral Flow Computation}
Interpolate boundary conditions by a phase $\varphi$,
\begin{equation}
    \psi(\theta+2\pi)=e^{i\varphi}\psi(\theta), \qquad \varphi\in[0,2\pi],
\end{equation}
which gives the eigenvalues
\begin{equation}
    \lambda_n(\varphi)=n+\frac{\varphi}{2\pi}, \qquad n\in\mathbb{Z}.
\end{equation}
As $\varphi$ runs from $0$ to $2\pi$, exactly one eigenvalue crosses zero, so
$\mathrm{SF}(U_\Gamma)=1$ and $\Delta_\Gamma=2\pi$.

\begin{remark}
For unitary matching maps $U_\Gamma:S^1\to U(1)$, the spectral flow equals the winding
number of $\det U_\Gamma$, consistent with the $K^1(S^1)\cong\mathbb{Z}$ classification.
\end{remark}

\subsection{\texorpdfstring{Normalization of $\phitree$}{Normalization of phi_tree}}
The two physical boundary cycles contribute $2\pi+2\pi$, and the seam gluing adds
$\Delta_\Gamma=2\pi$, so the total curvature is $R_{\text{total}}=6\pi$. The stationarity
normalization $\varphi\cdot R_{\text{total}}=1$ therefore yields
\begin{equation}
    \phitree = \frac{1}{6\pi}.
\end{equation}

%=============================================================================
\section{\texorpdfstring{Derivation of the Coefficient 48 in $\deltatop$}{Derivation of the Coefficient 48 in delta_top}}
\label{app:48deriv}
%=============================================================================

This appendix explains the integer coefficient in $\deltatop = 48\cthree^4$.
The key point is that the canonical quantity is
\begin{equation}
    \deltatop = \frac{3}{256\pi^4},
\end{equation}
and the appearance of ``48'' is simply the same number expressed in units of $\cthree^4$.

\subsection{The Physical Setup}

The topological correction $\deltatop$ arises from the leading gauge-invariant correction 
to the tree-level geometric scale $\phitree$. The requirements are:
\begin{enumerate}[leftmargin=1.5em, topsep=2pt, itemsep=1pt]
    \item \textbf{Gauge invariance:} The correction must be built from $\cthree$ only 
        (no explicit $\alpha$ dependence at this order).
    \item \textbf{Dimensionlessness:} Since $\cthree$ is dimensionless, the correction 
        must be a pure power $\cthree^n$ with an integer coefficient.
    \item \textbf{Discreteness:} Topological invariants take discrete values; the 
        coefficient must be a rational number with simple structure.
\end{enumerate}

\subsection{\texorpdfstring{Power Counting: Why $\cthree^4$?}{Power Counting: Why c3^4?}}

The effective action on the M\"obius fiber includes the Chern--Simons term 
$\cthree\, a F\tilde{F}$. The leading topological correction comes from the 
one-loop determinant of the Dirac operator coupled to this background.

\textbf{Dimensional analysis:} The Chern--Simons coupling enters as $\cthree$. The 
Euler characteristic and Pontryagin invariants contribute at order $\cthree^2$ each. 
The product yields:
\begin{equation}
    \deltatop \sim \cthree^2 \times \cthree^2 = \cthree^4
\end{equation}

\textbf{No $\cthree^2$ term:} A correction of order $\cthree^2$ would require a 
pure Euler term, but this is already absorbed into $\phitree$ via Gauss--Bonnet. 
The next allowed correction is $\cthree^4$.

\subsection{Spin-Lifted Deficit Factor}
The non-orientable identification acts as a $\mathbb{Z}_2$ operation on the base,
but in the spin bundle the same element lifts to order $4$ because a $2\pi$ rotation
acts as $-1$ on spinors and only a $4\pi$ rotation is trivial. This produces a
universal deficit factor
\begin{equation}
    1-\frac{1}{4}=\frac{3}{4},
\end{equation}
per minimal sector.

\subsection{Sector Count from Cut Geometry}
Cutting along $\Gamma$ yields boundary components $(C_1,C_2,\Gamma_+,\Gamma_-)$.
The four local corner patches
$(C_1,\Gamma_+)$, $(C_1,\Gamma_-)$, $(C_2,\Gamma_+)$, $(C_2,\Gamma_-)$
constitute the fundamental domain. Summing the spin-lifted deficit factor over these sectors yields
\begin{equation}
    \sum_{p=1}^{F}\left(1-\frac{1}{4}\right)=4\cdot\frac{3}{4}=3.
\end{equation}
In TFPT normalization this yields the canonical form $\deltatop = 3/(256\pi^4)$.

\subsection{\texorpdfstring{Conversion to the $\cthree$ normalization}{Conversion to the c3 normalization}}

Using $\cthree = 1/(8\pi)$:
\begin{equation}
    48\cthree^4 = 48 \times \frac{1}{(8\pi)^4} 
    = \frac{48}{4096\pi^4}
    = \frac{3}{256\pi^4}
\end{equation}
so $\deltatop = 3/(256\pi^4)$ and $\deltatop = 48\cthree^4$ are exactly equivalent.

\subsection{Status note}
The spin-lift argument fixes the $1-1/4$ deficit factor robustly, and the sector count
is fixed by seam symmetry. The seam normalization is provided by the $\eta$-gluing
derivation in Appendix~\ref{app:seam}.

%=============================================================================
\section{Flavor Monodromy and Cross-Ratio}
\label{app:flavor}
%=============================================================================

This appendix derives the M\"obius structure of flavor relations from boundary monodromy.

\subsection{Projective Invariance on Boundaries}

On a conformally flat boundary, the metric is locally:
\begin{equation}
    ds^2 = e^{2\omega(x)} \delta_{ij} dx^i dx^j
\end{equation}
Holomorphic automorphisms (Wilson lines of flat connections) act as M\"obius 
transformations:
\begin{equation}
    z \mapsto \frac{az + b}{cz + d}, \quad ad - bc = 1
\end{equation}

\begin{remark}
\textbf{Context:} $\mathrm{SL}(2,\mathbb{C})$ double-covers the proper orthochronous Lorentz group $\mathrm{SO}^+(3,1)$, 
and its projectivization $\mathrm{PSL}(2,\mathbb{C})$ acts by M\"obius transformations on the Riemann sphere. 
This links the M\"obius structure to standard Lorentz symmetry rather than an ad hoc choice
(see also the compact 4+2 note \cite{fourPlusTwoNote2025}).
\end{remark}

\subsection{Cross-Ratio as Unique Invariant}

Given four marked points $\{z_1, z_2, z_3, z_4\}$ on the boundary, the unique 
projective invariant is the cross-ratio:
\begin{equation}
    [z_1, z_2; z_3, z_4] = \frac{(z_1 - z_3)(z_2 - z_4)}{(z_1 - z_4)(z_2 - z_3)}
\end{equation}

\textbf{Consequence:} Any physical quantity depending only on boundary data must be 
expressible as a function of cross-ratios, hence as a M\"obius map.

\subsection{Cusp Selection}

The cusps $y \in \{1, 1/3, 2/3\}$ arise from:
\begin{enumerate}[leftmargin=1.5em, topsep=2pt, itemsep=1pt]
    \item \textbf{Rationality:} Quantized holonomies $\Rightarrow$ rational cusps.
    \item \textbf{GUT compatibility:} SU(5) hypercharge assignments give 
        $Y = 1/3$ (quarks), $Y = 1$ (leptons), etc.
    \item \textbf{Uniqueness:} These are the only rational values that simultaneously 
        separate the three SM sectors (up, down, lepton) and maintain GUT normalization.
\end{enumerate}

\subsection{\texorpdfstring{Derivation of $\delta_\star$}{Derivation of delta_star}}

The flavor phase is a first-order closure:
\begin{equation}
    \delta_\star = \delta_0 + \delta_1 \cdot \phiz + \mathcal{O}(\phiz^2)
\end{equation}

\textbf{Leading term:} $\delta_0 = 3/5$ from the abelian trace (same as $\bone$ numerator 
divided by 41/10's denominator structure).

\textbf{Correction:} $\delta_1 = 1/6$ from the M\"obius length scale. The factor $1/6$ 
matches the boundary normalization ($6\pi$ total curvature).

\textbf{Result:}
\begin{equation}
    \delta_\star = \frac{3}{5} + \frac{\phiz}{6} = 0.6 + 0.00886 = 0.6089
\end{equation}

%=============================================================================
\section{Reproducibility}
\label{app:repro}
%=============================================================================

\textbf{Replication package:} All numerical results can be reproduced using the 
TFPT verification suite (v5).

\textbf{Contents:}
\begin{itemize}[leftmargin=1.5em, topsep=2pt, itemsep=1pt]
    \item Python modules for each calculation domain
    \item JSON files with all constants and derived quantities
    \item Automated test suite with pass/fail status
    \item Plot generation scripts
    \item Verification suite modules include: \texttt{effective\_action\_r2}, \texttt{bounce\_perturbations} (see \codepath{tfpt-suite/} and the generated \codepath{tfpt-suite/tfpt-test-results.pdf})
\end{itemize}

\textbf{Key output files:}
\begin{itemize}[leftmargin=1.5em, topsep=2pt, itemsep=1pt]
    \item \texttt{tfpt\_complete\_analysis.json}: Master results file
    \item \texttt{e8\_chain\_analysis.json}: Complete E$_8$ chain data
    \item \texttt{gauge\_couplings.csv}: Two-loop RGE output (via PyR@TE)
\end{itemize}

\textbf{Execution:} \texttt{python -m tfpt\_verification5.run\_all}

\textbf{Verification:} The core modules must report ``success'' status; the
$R^2$ effective-action and bounce-perturbation modules are implemented in the shipped suite and included in the generated test report.

%=============================================================================
\section{Prediction Ledger and Falsification}
\label{app:prediction-ledger}
%=============================================================================

\textbf{Purpose:} Provide a single, reviewer-friendly list of predictions and candidates that is tied to the same JSON sources used by the verification suite, to prevent number drift between Abstract, theorems, tables, and appendices.

\begingroup
\footnotesize
\sisetup{round-mode=figures, round-precision=6}
\renewcommand{\arraystretch}{1.15}

% AUTO-GENERATED by tfpt_suite/export_prediction_ledger.py — inlined for standalone-paper build.
\begin{longtable}{@{}p{0.15\textwidth}p{0.27\textwidth}rrrp{0.10\textwidth}p{0.18\textwidth}@{}}
\caption{Prediction ledger (auto-generated from TFPT suite outputs).}\label{tab:prediction-ledger}\\
\toprule
Observable & Formula / dependencies & Pred. & Ref. & $z$ & Status & Suite evidence\\
\midrule
\endfirsthead
\toprule
Observable & Formula / dependencies & Pred. & Ref. & $z$ & Status & Suite evidence\\
\midrule
\endhead
\midrule
\multicolumn{7}{r}{\small Continued on next page.}\\
\midrule
\endfoot
\bottomrule
\endlastfoot
\multicolumn{7}{@{}l}{\textbf{Candidates / open interfaces}}\\
\(\Omega_b\) & \(\Omega_b=(4\pi-1)\beta_{\rm rad}\) under explicit sector-counting assumptions [Candidate] & \num{0.048940662665450112200545653760944651480054952073107413870556556782266783965399009} & \num{0.049301692328524438476424754994611856173233055557122787616860658651215012327847394} $\pm$ \num{0.00085681142747140931509835404112378635706093499964689237552400217734754363189007356} & \num{-0.421364} & C (conditional) & \texttt{omega\_b\_conjecture\_scan}\\ \codepath{out/omega_b_conjecture_scan/results.json}\\
\(\rho_\Lambda\) (normalization audit) & \(\phi_\*\!=n\,e^{-\alpha^{-1}(0)/2}\), \(\rho_\Lambda=(\bar M_{\rm Pl}\phi_\*)^4\); best \(n=1/2\) with log10 mismatch [Candidate] & \num{2.060247122268795528691070560351946768145020625187219242978373038664026793360454e-47} & \num{2.514176465474032e-47} & --- & C (PASS audit) & \texttt{dark\_energy\_exponential\_audit}\\ \codepath{theoryv3/out/dark_energy_exponential_audit/results.json}\\
\(\rho_\Lambda\) (torsion condensate closure) & Gap-equation completion with discrete \(n\) (spectral flow); compare to cosmology target [Candidate] & \num{4.1204942445375910573821411207038935362900412503744384859567460773280535867209079e-47} & \num{2.514176465474032e-47} & \num{2.1403530198484351956856249171986807549266690574189652264117407874776932146992054} & C (FAIL gate) & \texttt{torsion\_condensate}\\ \codepath{out_physics/torsion_condensate/results.json}\\
\multicolumn{7}{@{}l}{\textbf{Comparison layer}}\\
\(\mathrm{BR}(K^+\to\pi^+\nu\bar{\nu})\) & Rare kaon \(K\to\pi\nu\bar\nu\): \(\kappa_+\!\left[(\mathrm{Im}\,\lambda_t/\lambda^5\,X_t)^2+(\mathrm{Re}\,\lambda_c/\lambda\,P_c+\mathrm{Re}\,\lambda_t/\lambda^5\,X_t)^2\right]\) [Crosslink] & \num{8.221735267396824e-11} & \num{8.6e-11} $\pm$ \num{4.2e-12} & \num{-0.900630} & comparison & \texttt{kaon\_rare\_decay\_audit}\\ \codepath{out/kaon_rare_decay_audit/results.json}\\
\(\mathrm{BR}(K_L\to\pi^0\nu\bar{\nu})\) & Rare kaon \(K_L\to\pi^0\nu\bar\nu\): \(\kappa_L\,(\mathrm{Im}\,\lambda_t/\lambda^5\,X_t)^2\) (no direct measurement yet) [Crosslink] & \num{2.654823417968658e-11} & --- & --- & comparison & \texttt{kaon\_rare\_decay\_audit}\\ \codepath{out/kaon_rare_decay_audit/results.json}\\
\(\overline{\alpha}^{(5)}(M_Z)^{-1}\) & QED running + matching (declared policy) from \(\alpha(0)\) [Policy] & \num{127.94051870737564248749202130667255420059119636590364698278411016186034311550876} & \num{127.93} $\pm$ \num{0.008} & \num{1.314838421955310936502663334069275073899545737955872848013770232542889438594741} & comparison & \texttt{global\_consistency\_test}\\ \codepath{out/global_consistency_test/results.json}\\
\multicolumn{7}{@{}l}{\textbf{Conditional dynamics}}\\
\(A_s\) & Starobinsky \(R^2\) completion: \(A_s\simeq \frac{N^2}{24\pi^2}(M/\bar M_{\rm Pl})^2\) [Conditional] & \num{0.000000002090191364329459982659058205552757672085364626939873929875219205466615549540939} & \num{0.0000000020989031673} $\pm$ \num{0.00000000002938464434} & \num{-0.29647467805764898160213003439885527348994189187882714815192466383904684207796011} & conditional & \texttt{global\_consistency\_test}\\ \codepath{out/global_consistency_test/results.json}\\
\(\Omega_a h^2\) & Axion relic accounting (scenario-dependent) [Conditional] & \num{0.12275084115028702} & --- & --- & conditional & \texttt{axion\_dm\_audit}\\ \codepath{theoryv3/out/axion_dm_audit/results.json}\\
\(\nu\) [GHz] & \(\nu_{\rm GHz}\approx 0.24179893\, m_a[\mu\mathrm{eV}]\) [Conditional] & \num{15.76402318010866} & --- & --- & conditional & \texttt{axion\_dm\_audit}\\ \codepath{theoryv3/out/axion_dm_audit/results.json}\\
\(f_a\) [GeV] & PQ cascade block (axion claim) [Conditional] & \num{88639886850.34554} & --- & --- & conditional & \texttt{axion\_dm\_audit}\\ \codepath{theoryv3/out/axion_dm_audit/results.json}\\
\(m_a\) [\(\mu\)eV] & PQ cascade block (axion claim) [Conditional] & \num{65.1947598780055} & --- & --- & conditional & \texttt{axion\_dm\_audit}\\ \codepath{theoryv3/out/axion_dm_audit/results.json}\\
\(n_s\) & Starobinsky \(R^2\) completion: \(n_s=1-\frac{2}{N}\) (benchmark \(N=56\)) [Conditional] & \num{0.96428571428571428571428571428571428571428571428571428571428571428571428571428571} & \num{0.9649} $\pm$ \num{0.0042} & \num{-0.14625850340136054421768707482993197278911564625850340136054421768707482993197288} & conditional & \texttt{global\_consistency\_test}\\ \codepath{out/global_consistency_test/results.json}\\
\(r\) & Starobinsky \(R^2\) completion: \(r=\frac{12}{N^2}\) (benchmark \(N=56\)); compared to an upper bound [Conditional] & \num{0.003826530612244897959183673469387755102040816326530612244897959183673469387755102} & < 0.036 & --- & conditional & \texttt{global\_consistency\_test}\\ \codepath{out/global_consistency_test/results.json}\\
\multicolumn{7}{@{}l}{\textbf{Core (parameter-free)}}\\
\(\alpha^{-1}(0)\) & \(\mathrm{CFE} + \varphi_0(\alpha)\) with \(\delta_2=(g/4)\delta_{\mathrm{top}}^2\) and \(g=5\) [Kernel+Grammar] & \num{137.03599921615843479026171751711397398605568500258186146747743646328176776397918} & \num{137.035999177} $\pm$ \num{0.000000021} & \num{1.8646873709648436912911416183836040477419934032132112601562746554275798808614493} & P & \texttt{defect\_partition\_g5\_audit}\\ \codepath{theoryv3/out/defect_partition_g5_audit/results.json}\\
\(\beta_{\mathrm{deg}}\) & \(\beta_{\mathrm{deg}}=\frac{180}{\pi}\frac{\varphi_0}{4\pi}\) [Kernel] & \num{0.24243503090092952849243152047951445246859512239848881592836053961070607871616298} & \num{0.35} $\pm$ \num{0.14} & \num{-0.7683212078505033679112034251463253395100348400107941719402818599235280091702644} & P & \texttt{global\_consistency\_test}\\ \codepath{out/global_consistency_test/results.json}\\
\(\lambda\equiv |V_{us}|\) & \(\lambda=\sqrt{\varphi_0(1-\varphi_0)}\) [Kernel] & \num{0.22437623688472177311209717118786381569812652075494333477041354690545776148662783} & \num{0.22431} $\pm$ \num{0.00085} & \num{0.07792574673149777893784845631037140956061265287451149460417282995030763132686246} & P & \texttt{global\_consistency\_test}\\ \codepath{out/global_consistency_test/results.json}\\
\(\sin^2\theta_{12}\) & TM1 sum rule from \(\sin^2\theta_{13}\) [Kernel+Grammar] & \num{0.31792771} & \num{0.307} $\pm$ \num{0.012} & \num{0.910642} & P & \texttt{pmns\_z3\_breaking}\\ \codepath{out/pmns_z3_breaking/results.json}\\
\(\sin^2\theta_{13}\) & Identity + TM1 [Kernel+Grammar] & \num{0.02310844} & \num{0.02224} $\pm$ \num{0.00057} & \num{1.523570} & P & \texttt{pmns\_z3\_breaking}\\ \codepath{out/pmns_z3_breaking/results.json}\\
\(g_{a\gamma\gamma}\) & \(g_{a\gamma\gamma}=-4c_3=-\frac{1}{2\pi}\) [Kernel] & \num{-0.15915494309189535} & --- & --- & P & \texttt{(algebraic)}\newline \codepath{paper kernel (no suite output file)}\\
\end{longtable}
\endgroup

\textbf{Falsification guide (high-level):}
\begin{itemize}[leftmargin=1.5em, topsep=2pt, itemsep=1pt]
    \item \textbf{Core (parameter-free):} Any robust shift of CODATA $\alpha^{-1}(0)$ or a decisive inconsistency in the \texttt{defect\_partition\_g5\_audit} closure would directly falsify the TFPT kernel.
    \item \textbf{Comparison layer:} A stable discrepancy in $\overline{\alpha}^{(5)}(M_Z)^{-1}$ after adopting the declared running/matching policy falsifies the paper-to-observable bridge (not the kernel itself).
    \item \textbf{Conditional dynamics:} Failures here indicate that the adopted completion/policy layer (e.g.\ $R^2$ benchmark choices, relic accounting scenario) is not closed or not correct, and must be treated as conditional rather than unconditional prediction.
    \item \textbf{Candidates / open interfaces:} These rows are explicitly not ToE-closed; a FAIL gate is reported as such and should not be framed as a resolved prediction.
\end{itemize}

%=============================================================================
